\section{Metodologia}\label{sec:sv_metodo}

\textbf{Pré-processamento.} Os cinco atributos categóricos viram colunas 0/1, uma por categoria (\emph{one-hot}), e os cinco numéricos são padronizados (subtrai-se a média e divide-se pelo desvio padrão). São \svColunas{} colunas no total. Codificação, padronização e modelo ficam num único \emph{pipeline} do scikit-learn~\cite{pedregosa:11}, de modo que a média e o desvio vêm só do treino de cada divisão, e de cada parte da validação cruzada.

\textbf{Protocolo.} Os dados são divididos em treino (80\%, \svTreino{} pessoas) e teste (20\%, \svTeste, das quais \svPositivosTeste{} acima de 50 mil), mantendo a proporção das classes, e o procedimento é repetido com 10 seeds (0 a 9), que mudam a divisão, as partes da validação cruzada e a semente dos modelos. Quando um modelo tem hiperparâmetros, uma busca em grade com validação cruzada estratificada de 5 partes, só no treino, escolhe a combinação de maior acurácia, e o modelo é treinado de novo no treino inteiro. No teste, reportam-se a acurácia, o F1 macro (média do F1 das duas classes), o F1 da classe acima de 50 mil, a AUC e os falsos positivos e negativos, com a classe acima de 50 mil como positiva. A AUC usa a probabilidade prevista ou, na SVM, a distância à fronteira: mede se o modelo ordena bem as pessoas, independentemente do corte.

\textbf{Modelos.} As referências são a classe majoritária; a regra de um corte, uma árvore de profundidade 1 que escolhe sozinha, no treino, o melhor atributo e o melhor limiar; a regressão logística; a árvore de decisão com busca em grade; e a Random Forest com 300 árvores. A SVM aparece em quatro versões: a SVM linear, pelo LIBLINEAR (\texttt{LinearSVC}); a SVM com kernel RBF e a com kernel polinomial ($c_0 = 1$), pelo LIBSVM (\texttt{SVC}); e a SVM RBF com pesos por classe (\texttt{class\_weight="balanced"}, que faz o erro na classe menor custar cerca de 3 vezes mais), com os $C$ e $\gamma$ escolhidos para a RBF em cada seed. A rede neural é um MLP do scikit-learn: camadas ocultas com ativação ReLU, um neurônio de saída com sigmoide, perda de entropia cruzada e otimizador Adam, com parada antecipada (10\% do treino ficam de fora como validação, e o treino para depois de 30 épocas sem melhora). A busca da rede cobre a arquitetura (8, 16, 32 ou 64 neurônios numa camada; 32 e 16 ou 64 e 32 em duas), a regularização L2 e a taxa de aprendizado. As grades da SVM saíram de um estudo numa única divisão (a seed 0), apresentado na Seção~\ref{sec:sv_estudo}; a Tabela~\ref{tab:sv_grades} reúne todas.

\begin{table}[ht]
\centering
\caption{Grades da busca (número de combinações após os dois-pontos) e o que a busca escolheu nas 10 seeds.}
\label{tab:sv_grades}
\footnotesize
\begin{tabular}{l>{\raggedright\arraybackslash}p{5.6cm}>{\raggedright\arraybackslash}p{4.2cm}}
\toprule
Modelo & Valores testados & Escolha nas 10 seeds \\
\midrule
Árvore de decisão & critério (Gini, entropia) $\times$ profundidade (3, 5, 8, 12, sem limite) $\times$ mínimo por folha (1, 5, 20): 30 & 8 combinações diferentes; profundidade de 8 para cima ou sem limite \\
SVM linear & C $\in$ \{0,001; 0,01; 0,1; 1; 10\}: 5 & C = 0,1 em 7 \\
SVM RBF & C $\in$ \{0,1; 1; 10; 100\} $\times$ $\gamma$ $\in$ \{0,01; 0,03; 0,1; 0,3\}: 16 & C = 1 em 9, com $\gamma$ = 0,3 (5) ou 0,1 (4); C = 10 e $\gamma$ = 0,03 em 1 \\
SVM polinomial & grau $\in$ \{2; 3\} $\times$ C $\in$ \{0,1; 1; 10\}, coef0 = 1: 6 & grau 2 e C = 1 em 9 \\
Rede neural & 6 arquiteturas $\times$ $\alpha$ (10$^{-4}$, 10$^{-2}$) $\times$ taxa (10$^{-3}$, 10$^{-2}$): 24 & 10 combinações diferentes \\
\bottomrule
\end{tabular}

\end{table}

\textbf{Tempo de execução.} Para comparar algoritmos, e não o número de núcleos que cada um usa, tudo roda num único núcleo (\texttt{n\_jobs=1} e as bibliotecas de álgebra linear limitadas a uma \emph{thread}), com uma rodada de aquecimento descartada. Nos modelos com busca, o tempo da busca é medido à parte do treino final. O resumo usa a mediana das 10 seeds. A máquina é um AMD Ryzen 5 9600X (6 núcleos) com 31 GB de memória, Windows 11 e scikit-learn 1.9.1.
